\import{surreal/embedding.tex}
\import{surreal/sign-cut.tex}

\section{Dyadic surreals}

\subsection{Negative powers of two}

For a natural number $n$, the sign sequence $\surrealEpsilon{n}$ consists of
one positive sign followed by $n$ negative signs.  It represents $2^{-n}$.

\begin{definition}\label{surreal_epsilon}
    $\surrealEpsilon{n}=
        \{(\emptyset,\suc{\emptyset})\}
        \union
        \{(\suc{k},\emptyset)\mid k\in n\}$.
\end{definition}

\begin{proposition}\label{surreal_epsilon_iff}
    $z\in\surrealEpsilon{n}$ iff
        $z=(\emptyset,\suc{\emptyset})$ or
        there exists $k\in n$ such that $z=(\suc{k},\emptyset)$.
\end{proposition}
\begin{proof}
    \begin{align*}
        z\in\surrealEpsilon{n}
            &\iff
                z\in\{(\emptyset,\suc{\emptyset})\}\lor
                z\in\{(\suc{k},\emptyset)\mid k\in n\}
                \explanation{by \auto{surreal_epsilon,union_iff}}
            \\
            &\iff
                z=(\emptyset,\suc{\emptyset})\lor
                (\exists k\in n.\ z=(\suc{k},\emptyset))
                \explanation{by \auto{singleton_iff}}
    \end{align*}
    Follows by assumption.
\end{proof}

\begin{proposition}\label{surreal_epsilon_zero}
    $\surrealEpsilon{\emptyset}=\surrealOne$.
\end{proposition}
\begin{proof}
    Show that every element of $\surrealEpsilon{\emptyset}$ is an
        element of $\surrealOne$.
    \begin{subproof}
        Fix $z\in\surrealEpsilon{\emptyset}$.
        Hence $z=(\emptyset,\suc{\emptyset})$
            by \auto{surreal_epsilon_iff,emptyset}.
        Follows by
            \auto{ordinal_surreal,times_tuple_intro,suc_intro_self,singleton_iff}.
    \end{subproof}
    Show that every element of $\surrealOne$ is an element of
        $\surrealEpsilon{\emptyset}$.
    \begin{subproof}
        Fix $z\in\surrealOne$.
        Take $a,b$ such that $z=(a,b)$ and
            $a\in\suc{\emptyset}$ and
            $b\in\{\suc{\emptyset}\}$
            by \auto{ordinal_surreal,times_elem_is_tuple}.
        Hence $a=\emptyset$ and $b=\suc{\emptyset}$
            by \obv{suc_iff,emptyset,singleton_iff}.
        Follows by \auto{surreal_epsilon_iff}.
    \end{subproof}
    Follows by \obv{setext}.
\end{proof}

\begin{lemma}\label{surreal_double_injective}
    Suppose $\surreal{x}$ and $\surreal{y}$.
    Suppose $\surrealAdd{x}{x}=\surrealAdd{y}{y}$.
    Then $x=y$.
\end{lemma}
\begin{proof}
    $\surrealLt{x}{y}$ or $x=y$ or $\surrealLt{y}{x}$
        by \obv{surreal_lt_trichotomy}.
    \begin{byCase}
        \caseOf{$x=y$.}
            Follows.
        \caseOf{$x\neq y$.}
            \begin{byCase}
                \caseOf{$\surrealLt{x}{y}$.}
                    $\surrealLt{\surrealAdd{x}{x}}
                        {\surrealAdd{y}{y}}$
                        by \obv{surreal_add_lt_both}.
                    Contradiction by
                        \obv{surreal_lt_irreflexive,surreal_add_is_surreal}.
                \caseOf{not $\surrealLt{x}{y}$.}
                    Hence $\surrealLt{y}{x}$.
                    $\surrealLt{\surrealAdd{y}{y}}
                        {\surrealAdd{x}{x}}$
                        by \obv{surreal_add_lt_both}.
                    Contradiction by
                        \obv{surreal_lt_irreflexive,surreal_add_is_surreal}.
            \end{byCase}
    \end{byCase}
\end{proof}

\begin{proposition}\label{surreal_epsilon_function}
    Let $n$ be a natural number.
    Then $\surrealEpsilon{n}\in
        \funs{\suc{n}}{\surrealSigns}$.
\end{proposition}
\begin{proof}
    $n$ and $\suc{n}$ are ordinals
        by \obv{natural_number_is_ordinal,suc_ordinal}.
    Show that $\surrealEpsilon{n}$ is a relation.
    \begin{subproof}
        Show that for every $z\in\surrealEpsilon{n}$ there exist $a,b$
            such that $z=(a,b)$.
        \begin{subproof}
            Fix $z\in\surrealEpsilon{n}$.
            Follows by \auto{surreal_epsilon_iff}.
        \end{subproof}
        Follows by \obv{relation}.
    \end{subproof}
    Show that $\surrealEpsilon{n}$ is right-unique.
    \begin{subproof}
        Show that for all $a,b,c$ such that
            $(a,b),(a,c)\in\surrealEpsilon{n}$ we have $b=c$.
        \begin{subproof}
            Fix $a,b,c$.
            Assume $(a,b),(a,c)\in\surrealEpsilon{n}$.
            Follows by \auto{surreal_epsilon_iff,pair_eq_iff,suc_neq_emptyset}.
        \end{subproof}
        Follows by \obv{rightunique}.
    \end{subproof}

    Show that $\dom{\surrealEpsilon{n}}\subseteq\suc{n}$.
    \begin{subproof}
        Show that for every $a\in\dom{\surrealEpsilon{n}}$ we have
            $a\in\suc{n}$.
        \begin{subproof}
            Fix $a\in\dom{\surrealEpsilon{n}}$.
            Take $b$ such that $(a,b)\in\surrealEpsilon{n}$
                by \obv{dom_iff}.
            $a=\emptyset$ or there exists $k\in n$ such that
                $a=\suc{k}$
                by \auto{surreal_epsilon_iff,pair_eq_iff}.
            \begin{byCase}
                \caseOf{$a=\emptyset$.}
                    $\emptyset\in n$ or $\emptyset=n$
                        by \obv{ordinal_empty_or_emptyset_elem}.
                    Follows by \auto{suc_iff}.
                \caseOf{$a\neq\emptyset$.}
                    Take $k\in n$ such that $a=\suc{k}$.
                    $k$ is an ordinal by \obv{elem_of_ordinal_is_ordinal}.
                    $\suc{k}\subseteq n$ by \obv{ordinal_suc_subseteq}.
                    Hence $\suc{k}\in n$ or $\suc{k}=n$
                        by \obv{ordinal_subseteq_implies_elem_or_eq,suc_ordinal}.
                    Follows by \auto{suc_iff}.
            \end{byCase}
        \end{subproof}
        Follows by \obv{subseteq}.
    \end{subproof}
    Show that $\suc{n}\subseteq\dom{\surrealEpsilon{n}}$.
    \begin{subproof}
        Show that for every $a\in\suc{n}$ we have
            $a\in\dom{\surrealEpsilon{n}}$.
        \begin{subproof}
            Fix $a\in\suc{n}$.
            $\suc{n}$ is a natural number by \obv{omega_successor}.
            Hence $a$ is a natural number
                by \obv{naturals_is_transitive_set,transitiveset_iff_subseteq,subseteq}.
            $a=\emptyset$ or there exists a natural number $k$ such that
                $a=\suc{k}$
                by \obv{natural_zero_or_successor}.
            \begin{byCase}
                \caseOf{$a=\emptyset$.}
                    $(a,\suc{\emptyset})\in\surrealEpsilon{n}$
                        by \auto{surreal_epsilon_iff}.
                    Follows by \obv{dom_intro}.
                \caseOf{$a\neq\emptyset$.}
                    Take $k$ such that $k$ is a natural number and
                        $a=\suc{k}$.
                    $a\in n$ or $a=n$ by \auto{suc_iff}.
                    $k\in a$ by \obv{suc_intro_self}.
                    Hence $k\in n$
                        by \obv{ordinal_transitivity}.
                    $(a,\emptyset)\in\surrealEpsilon{n}$
                        by \auto{surreal_epsilon_iff}.
                    Follows by \obv{dom_intro}.
            \end{byCase}
        \end{subproof}
        Follows by \obv{subseteq}.
    \end{subproof}
    Hence $\dom{\surrealEpsilon{n}}=\suc{n}$
        by \obv{subseteq_antisymmetric}.
    For every $a\in\dom{\surrealEpsilon{n}}$ we have
        $\apply{\surrealEpsilon{n}}{a}\in\surrealSigns$
        by \auto{function_apply_elim,surreal_epsilon_iff,pair_eq_iff,surreal_sign_members}.
    Thus $\surrealEpsilon{n}$ is a function to $\surrealSigns$.
    Follows by \obv{funs_intro}.
\end{proof}

\begin{proposition}\label{surreal_epsilon_at}
    Let $n$ be a natural number.
    Then $\surrealAt{\suc{n}}{\surrealEpsilon{n}}$.
\end{proposition}
\begin{proof}
    $\suc{n}$ is an ordinal
        by \obv{natural_number_is_ordinal,suc_ordinal}.
    Follows by \obv{surreal_at,surreal_epsilon_function}.
\end{proof}

\begin{corollary}\label{surreal_epsilon_is_surreal}
    Let $n$ be a natural number.
    Then $\surreal{\surrealEpsilon{n}}$ and
        $\birthday{\surrealEpsilon{n}}=\suc{n}$.
\end{corollary}
\begin{proof}
    Follows by \obv{surreal_epsilon_at,surreal_at_implies_surreal}.
\end{proof}

\begin{proposition}\label{surreal_epsilon_positive_value}
    Let $n$ be a natural number.
    Let $\alpha\in\suc{n}$.
    Then $\apply{\surrealEpsilon{n}}{\alpha}=\suc{\emptyset}$ iff
        $\alpha=\emptyset$.
\end{proposition}
\begin{proof}
    $\surrealEpsilon{n}$ is a function and
        $\alpha\in\dom{\surrealEpsilon{n}}$
        by \auto{surreal_epsilon_function,funs_is_function,funs_elim}.
    $(\alpha,\apply{\surrealEpsilon{n}}{\alpha})
        \in\surrealEpsilon{n}$
        by \obv{function_apply_elim}.
    $(\emptyset,\suc{\emptyset})\in\surrealEpsilon{n}$
        by \auto{surreal_epsilon_iff}.
    Follows by \auto{surreal_epsilon_iff,function_apply_intro,pair_eq_iff,suc_neq_emptyset}.
\end{proof}

\begin{proposition}\label{surreal_epsilon_negative_value}
    Let $n$ be a natural number.
    Let $\alpha\in\suc{n}$.
    Then $\apply{\surrealEpsilon{n}}{\alpha}=\emptyset$ iff
        there exists $k\in n$ such that $\alpha=\suc{k}$.
\end{proposition}
\begin{proof}
    $\surrealEpsilon{n}$ is a function and
        $\alpha\in\dom{\surrealEpsilon{n}}$
        by \auto{surreal_epsilon_function,funs_is_function,funs_elim}.
    $(\alpha,\apply{\surrealEpsilon{n}}{\alpha})
        \in\surrealEpsilon{n}$
        by \obv{function_apply_elim}.
    For every $k\in n$ we have
        $(\suc{k},\emptyset)\in\surrealEpsilon{n}$
        by \auto{surreal_epsilon_iff}.
    Follows by \auto{surreal_epsilon_iff,function_apply_intro,pair_eq_iff,suc_neq_emptyset}.
\end{proof}

\begin{proposition}\label{surreal_epsilon_restriction_empty}
    $\restrl{\surrealEpsilon{n}}{\emptyset}=\emptyset$.
\end{proposition}
\begin{proof}
    Follows by \auto{setext,restrl,emptyset}.
\end{proof}

\begin{proposition}\label{surreal_epsilon_restriction}
    Let $n$ be a natural number.
    Let $k\in n$.
    Then $\restrl{\surrealEpsilon{n}}{\suc{k}}
        =\surrealEpsilon{k}$.
\end{proposition}
\begin{proof}
    $k$ is a natural number
        by \obv{naturals_is_transitive_set,transitiveset_iff_subseteq,subseteq}.
    $n$ and $k$ are ordinals by \obv{natural_number_is_ordinal}.
    Show that every element of
        $\restrl{\surrealEpsilon{n}}{\suc{k}}$ is an element of
        $\surrealEpsilon{k}$.
    \begin{subproof}
        Fix $z\in\restrl{\surrealEpsilon{n}}{\suc{k}}$.
        Take $a,b$ such that $z=(a,b)$ and $a\in\suc{k}$ and
            $(a,b)\in\surrealEpsilon{n}$
            by \auto{restrl}.
        $z=(\emptyset,\suc{\emptyset})$ or
            there exists $j\in n$ such that
                $z=(\suc{j},\emptyset)$
            by \obv{surreal_epsilon_iff}.
        \begin{byCase}
            \caseOf{$z=(\emptyset,\suc{\emptyset})$.}
                Follows by \auto{surreal_epsilon_iff}.
            \caseOf{$z\neq(\emptyset,\suc{\emptyset})$.}
                Take $j\in n$ such that $z=(\suc{j},\emptyset)$.
                Hence $a=\suc{j}$ by \auto{pair_eq_iff}.
                $\suc{j}\in k$ or $\suc{j}=k$
                    by \auto{suc_iff}.
                $j\in\suc{j}$ by \obv{suc_intro_self}.
                Hence $j\in k$ by \obv{ordinal_transitivity}.
                Follows by \auto{surreal_epsilon_iff}.
        \end{byCase}
    \end{subproof}
    Show that every element of $\surrealEpsilon{k}$ is an element of
        $\restrl{\surrealEpsilon{n}}{\suc{k}}$.
    \begin{subproof}
        Fix $z\in\surrealEpsilon{k}$.
        $z=(\emptyset,\suc{\emptyset})$ or
            there exists $j\in k$ such that
                $z=(\suc{j},\emptyset)$
            by \obv{surreal_epsilon_iff}.
        \begin{byCase}
            \caseOf{$z=(\emptyset,\suc{\emptyset})$.}
                $\emptyset\in k$ or $\emptyset=k$
                    by \obv{ordinal_empty_or_emptyset_elem}.
                Hence $\emptyset\in\suc{k}$ by \auto{suc_iff}.
                $z\in\surrealEpsilon{n}$
                    by \auto{surreal_epsilon_iff}.
                Follows by \auto{restrl}.
            \caseOf{$z\neq(\emptyset,\suc{\emptyset})$.}
                Take $j\in k$ such that $z=(\suc{j},\emptyset)$.
                $j\in n$ by \obv{ordinal_transitivity}.
                Hence $z\in\surrealEpsilon{n}$
                    by \auto{surreal_epsilon_iff}.
                $j$ is an ordinal by \obv{elem_of_ordinal_is_ordinal}.
                $\suc{j}\subseteq k$ by \obv{ordinal_suc_subseteq}.
                Hence $\suc{j}\in k$ or $\suc{j}=k$
                    by \obv{ordinal_subseteq_implies_elem_or_eq,suc_ordinal}.
                Hence $\suc{j}\in\suc{k}$ by \auto{suc_iff}.
                Follows by \auto{restrl}.
        \end{byCase}
    \end{subproof}
    Follows by \obv{setext}.
\end{proof}

\begin{proposition}\label{surreal_epsilon_image_iff}
    $x\in\{\surrealEpsilon{k}\mid k\in n\}$ iff
        there exists $k\in n$ such that $x=\surrealEpsilon{k}$.
\end{proposition}
\begin{proof}
    Follows by \auto{logical_replacement_iff}.
\end{proof}

\begin{proposition}\label{surreal_epsilon_sign_left}
    Let $n$ be a natural number.
    Then $\surrealSignLeft{\surrealEpsilon{n}}=\{\emptyset\}$.
\end{proposition}
\begin{proof}
    Show that every element of $\surrealSignLeft{\surrealEpsilon{n}}$
        is equal to $\emptyset$.
    \begin{subproof}
        Fix $l\in\surrealSignLeft{\surrealEpsilon{n}}$.
        Take $\alpha\in\birthday{\surrealEpsilon{n}}$ such that
            $\apply{\surrealEpsilon{n}}{\alpha}=\suc{\emptyset}$ and
            $l=\restrl{\surrealEpsilon{n}}{\alpha}$
            by \obv{surreal_sign_left_iff}.
        $\birthday{\surrealEpsilon{n}}=\suc{n}$
            by \obv{surreal_epsilon_is_surreal}.
        Hence $\alpha=\emptyset$
            by \obv{surreal_epsilon_positive_value}.
        Follows by \auto{surreal_epsilon_restriction_empty}.
    \end{subproof}
    $\emptyset\in\suc{n}$
        by \auto{natural_number_is_ordinal,ordinal_empty_or_emptyset_elem,suc_iff}.
    $\apply{\surrealEpsilon{n}}{\emptyset}=\suc{\emptyset}$
        by \auto{surreal_epsilon_positive_value}.
    $\restrl{\surrealEpsilon{n}}{\emptyset}=\emptyset$
        by \obv{surreal_epsilon_restriction_empty}.
    Hence $\emptyset\in\surrealSignLeft{\surrealEpsilon{n}}$
        by \auto{surreal_sign_left_iff,surreal_epsilon_is_surreal}.
    Follows by \auto{setext,singleton_iff}.
\end{proof}

\begin{proposition}\label{surreal_epsilon_sign_right}
    Let $n$ be a natural number.
    Then $\surrealSignRight{\surrealEpsilon{n}}=
        \{\surrealEpsilon{k}\mid k\in n\}$.
\end{proposition}
\begin{proof}
    Show that every element of $\surrealSignRight{\surrealEpsilon{n}}$
        is an element of $\{\surrealEpsilon{k}\mid k\in n\}$.
    \begin{subproof}
        Fix $r\in\surrealSignRight{\surrealEpsilon{n}}$.
        Take $\alpha\in\birthday{\surrealEpsilon{n}}$ such that
            $\apply{\surrealEpsilon{n}}{\alpha}=\emptyset$ and
            $r=\restrl{\surrealEpsilon{n}}{\alpha}$
            by \obv{surreal_sign_right_iff}.
        $\birthday{\surrealEpsilon{n}}=\suc{n}$
            by \obv{surreal_epsilon_is_surreal}.
        Take $k\in n$ such that $\alpha=\suc{k}$
            by \obv{surreal_epsilon_negative_value}.
        $r=\surrealEpsilon{k}$
            by \auto{surreal_epsilon_restriction}.
        Follows by \auto{surreal_epsilon_image_iff}.
    \end{subproof}
    Show that every element of $\{\surrealEpsilon{k}\mid k\in n\}$
        is an element of $\surrealSignRight{\surrealEpsilon{n}}$.
    \begin{subproof}
        Fix $r\in\{\surrealEpsilon{k}\mid k\in n\}$.
        Take $k\in n$ such that $r=\surrealEpsilon{k}$
            by \obv{surreal_epsilon_image_iff}.
        $k$ is a natural number
            by \obv{naturals_is_transitive_set,transitiveset_iff_subseteq,subseteq}.
        $n$ and $k$ are ordinals by \obv{natural_number_is_ordinal}.
        $\suc{k}\subseteq n$ by \obv{ordinal_suc_subseteq}.
        Hence $\suc{k}\in n$ or $\suc{k}=n$
            by \obv{ordinal_subseteq_implies_elem_or_eq,suc_ordinal}.
        Hence $\suc{k}\in\suc{n}$ by \auto{suc_iff}.
        $\birthday{\surrealEpsilon{n}}=\suc{n}$
            by \obv{surreal_epsilon_is_surreal}.
        $\apply{\surrealEpsilon{n}}{\suc{k}}=\emptyset$
            by \auto{surreal_epsilon_negative_value}.
        $\restrl{\surrealEpsilon{n}}{\suc{k}}
            =\surrealEpsilon{k}$
            by \obv{surreal_epsilon_restriction}.
        Follows by \auto{surreal_sign_right_iff}.
    \end{subproof}
    Follows by \obv{setext}.
\end{proof}

\begin{theorem}\label{surreal_epsilon_cut_equation}
    Let $n$ be a natural number.
    Then $\surrealEpsilon{n}=
        \surrealCut{\{\emptyset\}}
            {\{\surrealEpsilon{k}\mid k\in n\}}$.
\end{theorem}
\begin{proof}
    $\surreal{\surrealEpsilon{n}}$
        by \auto{surreal_epsilon_is_surreal}.
    $\surrealEpsilon{n}=\surrealCut
        {\surrealSignLeft{\surrealEpsilon{n}}}
        {\surrealSignRight{\surrealEpsilon{n}}}$
        by \obv{surreal_sign_cut_equation}.
    Follows by
        \obv{surreal_epsilon_sign_left,surreal_epsilon_sign_right}.
\end{proof}

\begin{proposition}\label{surreal_epsilon_positive}
    Let $n$ be a natural number.
    Then $\surrealLt{\surrealZero}{\surrealEpsilon{n}}$.
\end{proposition}
\begin{proof}
    $\surreal{\surrealEpsilon{n}}$ and
        $\birthday{\surrealEpsilon{n}}=\suc{n}$
        by \obv{surreal_epsilon_is_surreal}.
    $\birthday{\emptyset}=\emptyset$
        by \obv{empty_surreal}.
    $\emptyset\in\suc{n}$
        by \auto{natural_number_is_ordinal,ordinal_empty_or_emptyset_elem,suc_iff}.
    $\surrealEq{\emptyset}{\emptyset}{\surrealEpsilon{n}}$
        by \auto{surreal_eq_initial}.
    $\apply{\surrealEpsilon{n}}{\emptyset}=\suc{\emptyset}$
        by \auto{surreal_epsilon_positive_value}.
    Follows by \obv{surreal_lt,surreal_lt_left_short}.
\end{proof}

\begin{theorem}\label{surreal_epsilon_decreasing}
    Let $n$ be a natural number.
    Let $k\in n$.
    Then $\surrealLt{\surrealEpsilon{n}}{\surrealEpsilon{k}}$.
\end{theorem}
\begin{proof}
    $k$ is a natural number
        by \obv{naturals_is_transitive_set,transitiveset_iff_subseteq,subseteq}.
    $n$ and $k$ are ordinals by \obv{natural_number_is_ordinal}.
    $\suc{k}\subseteq n$ by \obv{ordinal_suc_subseteq}.
    Hence $\suc{k}\in n$ or $\suc{k}=n$
        by \obv{ordinal_subseteq_implies_elem_or_eq,suc_ordinal}.
    Hence $\suc{k}\in\suc{n}$ by \auto{suc_iff}.
    $\birthday{\surrealEpsilon{n}}=\suc{n}$ and
        $\birthday{\surrealEpsilon{k}}=\suc{k}$
        by \obv{surreal_epsilon_is_surreal}.
    $\surrealEq{\suc{k}}
        {\restrl{\surrealEpsilon{n}}{\suc{k}}}
        {\surrealEpsilon{n}}$
        by \obv{surreal_restriction_agrees,surreal_epsilon_is_surreal}.
    $\restrl{\surrealEpsilon{n}}{\suc{k}}
        =\surrealEpsilon{k}$
        by \obv{surreal_epsilon_restriction}.
    Hence $\surrealEq{\suc{k}}
        {\surrealEpsilon{n}}{\surrealEpsilon{k}}$
        by \obv{surreal_eq_initial_symmetric}.
    $\apply{\surrealEpsilon{n}}{\suc{k}}=\emptyset$
        by \auto{surreal_epsilon_negative_value}.
    Follows by \obv{surreal_lt,surreal_lt_right_short}.
\end{proof}

\begin{corollary}\label{surreal_epsilon_antitone}
    Let $m,n$ be natural numbers.
    Suppose $m\subseteq n$.
    Then $\surrealLe{\surrealEpsilon{n}}{\surrealEpsilon{m}}$.
\end{corollary}
\begin{proof}
    $m$ and $n$ are ordinals by \obv{natural_number_is_ordinal}.
    Hence $m\in n$ or $m=n$ by
        \obv{ordinal_subseteq_implies_elem_or_eq}.
    \begin{byCase}
        \caseOf{$m\in n$.}
            $\surrealLt{\surrealEpsilon{n}}{\surrealEpsilon{m}}$ by
                \obv{surreal_epsilon_decreasing}.
            Follows by
                \obv{surreal_le,surreal_epsilon_is_surreal}.
        \caseOf{$m\notin n$.}
            Hence $m=n$.
            Follows by
                \obv{surreal_le_reflexive,surreal_epsilon_is_surreal}.
    \end{byCase}
\end{proof}

\begin{theorem}\label{surreal_epsilon_halving}
    Let $n$ be an ordinal.
    Suppose $n$ is a natural number.
    Then
        $\surrealAdd{\surrealEpsilon{\suc{n}}}
            {\surrealEpsilon{\suc{n}}}=\surrealEpsilon{n}$.
\end{theorem}
\begin{proof}[Proof by transfinite induction]
    Let $x=\surrealEpsilon{\suc{n}}$.
    Let $A=\{\emptyset\}$.
    Let $B=\{\surrealEpsilon{k}\mid k\in\suc{n}\}$.
    Let $L=\surrealAddCutLeft{A}{A}{x}{x}$.
    Let $R=\surrealAddCutRight{B}{B}{x}{x}$.
    Let $M=\{\emptyset\}$.
    Let $N=\{\surrealEpsilon{k}\mid k\in n\}$.

    $\suc{n}$ is a natural number by \obv{omega_successor}.
    $\surreal{x}$ and $\surreal{\surrealEpsilon{n}}$
        by \auto{surreal_epsilon_is_surreal}.
    $\surrealLt{\surrealZero}{x}$
        by \obv{surreal_epsilon_positive}.
    $\surrealLt{x}{\surrealEpsilon{n}}$
        by \auto{surreal_epsilon_decreasing,suc_intro_self}.

    $\surrealCutPair{A}{B}$ and $x=\surrealCut{A}{B}$
        by \auto{surreal_sign_cut_pair,surreal_epsilon_sign_left,surreal_epsilon_sign_right,surreal_epsilon_cut_equation}.
    Hence $\surrealCutPair{L}{R}$
        by \auto{surreal_add_general_cut_pair}.
    \begin{align*}
        \surrealAdd{x}{x}
            &= \surrealAdd{\surrealCut{A}{B}}
                {\surrealCut{A}{B}} \\
            &= \surrealCut{
                \surrealAddCutLeft{A}{A}
                    {\surrealCut{A}{B}}{\surrealCut{A}{B}}}{
                \surrealAddCutRight{B}{B}
                    {\surrealCut{A}{B}}{\surrealCut{A}{B}}}
                \explanation{by \obv{surreal_add_cuts_equation}}
            \\
            &= \surrealCut{L}{R}
    \end{align*}
    $\surrealCutPair{M}{N}$ and
        $\surrealEpsilon{n}=\surrealCut{M}{N}$
        by \auto{surreal_sign_cut_pair,surreal_epsilon_sign_left,surreal_epsilon_sign_right,surreal_epsilon_cut_equation}.

    Show that for every $l\in L$ we have
        $\surrealLt{l}{\surrealCut{M}{N}}$.
    \begin{subproof}
        Fix $l\in L$.
        $(\exists a:\Set.\ a\in A\land l=\surrealAdd{a}{x})
            \lor
            (\exists a:\Set.\ a\in A\land l=\surrealAdd{x}{a})$
            by \obv{surreal_add_cut_left_iff}.
        Hence $l=x$
            by \auto{singleton_iff,surreal_add_zero_left,surreal_add_zero_right}.
    \end{subproof}

    Show that for every $r\in R$ we have
        $\surrealLt{\surrealCut{M}{N}}{r}$.
    \begin{subproof}
        Fix $r\in R$.
        $(\exists k:\Set.\ k\in\suc{n}\land
                r=\surrealAdd{\surrealEpsilon{k}}{x})
            \lor
            (\exists k:\Set.\ k\in\suc{n}\land
                r=\surrealAdd{x}{\surrealEpsilon{k}})$
            by \auto{surreal_add_cut_right_iff}.
        Take $k\in\suc{n}$ such that
            $r=\surrealAdd{\surrealEpsilon{k}}{x}$ or
            $r=\surrealAdd{x}{\surrealEpsilon{k}}$ by \auto.
        $k$ is a natural number
            by \obv{naturals_is_transitive_set,transitiveset}.
        $k\in n$ or $k=n$ by \auto{suc_iff}.
        Show that
            $\surrealLt{\surrealEpsilon{n}}
                {\surrealAdd{\surrealEpsilon{k}}{x}}$.
        \begin{subproof}
            \begin{byCase}
                \caseOf{$k=n$.}
                    $\surrealLt
                        {\surrealAdd{\surrealEpsilon{n}}{\surrealZero}}
                        {\surrealAdd{\surrealEpsilon{n}}{x}}$
                        by \obv{surreal_add_lt_right,empty_surreal}.
                    Follows by
                        \obv{surreal_add_zero_right,surreal_epsilon_is_surreal}.
                \caseOf{$k\neq n$.}
                    Hence $k\in n$.
                    $\surrealLt{\surrealEpsilon{n}}{\surrealEpsilon{k}}$
                        by \obv{surreal_epsilon_decreasing}.
                    $\surrealLt{\surrealEpsilon{k}}
                        {\surrealAdd{\surrealEpsilon{k}}{x}}$
                        by \auto{surreal_add_lt_right,surreal_add_zero_right,surreal_epsilon_is_surreal,empty_surreal}.
                    Follows by
                        \auto{surreal_lt_transitive,surreal_add_is_surreal,surreal_epsilon_is_surreal}.
            \end{byCase}
        \end{subproof}
        $\surrealAdd{\surrealEpsilon{k}}{x}
            =\surrealAdd{x}{\surrealEpsilon{k}}$
            by \auto{surreal_add_commutative,surreal_epsilon_is_surreal}.
    \end{subproof}

    Show that for every $m\in M$ we have
        $\surrealLt{m}{\surrealCut{L}{R}}$.
    \begin{subproof}
        Fix $m\in M$.
        Hence $m=\surrealZero$ by \obv{singleton_iff}.
        $\surrealLt
            {\surrealAdd{\surrealZero}{\surrealZero}}
            {\surrealAdd{x}{x}}$
            by \obv{surreal_add_lt_both,empty_surreal}.
        Follows by
            \obv{surreal_add_zero_left,empty_surreal}.
    \end{subproof}

    Show that for every $r\in N$ we have
        $\surrealLt{\surrealCut{L}{R}}{r}$.
    \begin{subproof}
        Fix $r\in N$.
        Take $k\in n$ such that $r=\surrealEpsilon{k}$
            by \obv{surreal_epsilon_image_iff}.
        $k$ is a natural number
            by \obv{naturals_is_transitive_set,transitiveset}.
        $\suc{k}$ is a natural number by \obv{omega_successor}.
        $\suc{k}\in\suc{n}$
            by \auto{natural_number_is_ordinal,ordinal_suc_subseteq,ordinal_subseteq_implies_elem_or_eq,suc_ordinal,suc_iff}.
        Hence $\surrealLt{x}{\surrealEpsilon{\suc{k}}}$
            by \auto{surreal_epsilon_decreasing}.
        Hence
            $\surrealLt{\surrealAdd{x}{x}}
                {\surrealAdd{\surrealEpsilon{\suc{k}}}
                    {\surrealEpsilon{\suc{k}}}}$
            by \auto{surreal_add_lt_both,surreal_epsilon_is_surreal}.
        $\surrealAdd{\surrealEpsilon{\suc{k}}}
            {\surrealEpsilon{\suc{k}}}=\surrealEpsilon{k}$.
    \end{subproof}
    $\surrealCut{L}{R}=\surrealCut{M}{N}$
        by \obv{surreal_cut_extensionality}.
\end{proof}

\begin{theorem}\label{surreal_epsilon_product}
    Let $m$ be a natural number.
    Let $n$ be an ordinal.
    Suppose $n$ is a natural number.
    Then $\surrealMul{\surrealEpsilon{m}}{\surrealEpsilon{n}}
        =\surrealEpsilon{m+n}$.
\end{theorem}
\begin{proof}[Proof by transfinite induction]
    $n=\emptyset$ or there exists a natural number $k$ such that
        $n=\suc{k}$ by \obv{natural_zero_or_successor}.
    \begin{byCase}
        \caseOf{$n=\emptyset$.}
            $m+n=m$ by \auto{num_addition_zero_right}.
            $\surrealEpsilon{n}=\surrealOne$
                by \obv{surreal_epsilon_zero}.
            Follows by
                \auto{surreal_mul_one_right,surreal_epsilon_is_surreal}.
        \caseOf{$n\neq\emptyset$.}
            Take $k$ such that $k$ is a natural number and
                $n=\suc{k}$.
            $k\in n$ by \obv{suc_intro_self}.
            $k$ is an ordinal by \obv{natural_number_is_ordinal}.
            $m+k$ is a natural number by \obv{num_addition_on_naturals}.
            $m+n$ is a natural number by \obv{num_addition_on_naturals}.
            $m+n=\suc{m+k}$
                by \auto{num_addition_successor_right}.
            $\surrealMul{\surrealEpsilon{m}}{\surrealEpsilon{k}}
                =\surrealEpsilon{m+k}$.
            $\surrealAdd{\surrealEpsilon{n}}{\surrealEpsilon{n}}
                =\surrealEpsilon{k}$
                by \auto{surreal_epsilon_halving}.
            Let $a=\surrealMul{\surrealEpsilon{m}}
                {\surrealEpsilon{n}}$.
            Let $b=\surrealEpsilon{m+n}$.
            $\surreal{\surrealEpsilon{m}}$ and
                $\surreal{\surrealEpsilon{n}}$
                by \auto{surreal_epsilon_is_surreal}.
            Hence $\surreal{a}$ by \obv{surreal_mul_is_surreal}.
            $\surreal{b}$ by \auto{surreal_epsilon_is_surreal}.
            $\surrealAdd{a}{a}
                =\surrealMul{\surrealEpsilon{m}}
                    {\surrealAdd{\surrealEpsilon{n}}{\surrealEpsilon{n}}}$
                by \auto{surreal_mul_distributive_left,surreal_epsilon_is_surreal}.
            Hence $\surrealAdd{a}{a}=\surrealEpsilon{m+k}$.
            $\surrealAdd{b}{b}=\surrealEpsilon{m+k}$
                by \auto{surreal_epsilon_halving,natural_number_is_ordinal}.
            Hence $\surrealAdd{a}{a}=\surrealAdd{b}{b}$.
            Follows by \obv{surreal_double_injective}.
    \end{byCase}
\end{proof}

\begin{theorem}\label{surreal_epsilon_integer_rescale}
    Let $k$ be a natural number.
    Let $z\in\surrealIntegers$.
    Let $n$ be an ordinal.
    Suppose $n$ is a natural number.
    Then there exists $w\in\surrealIntegers$ such that
        $\surrealMul{\surrealEpsilon{k}}{z}
            =\surrealMul{\surrealEpsilon{k+n}}{w}$.
\end{theorem}
\begin{proof}[Proof by transfinite induction]
    $n=\emptyset$ or there exists a natural number $j$ such that
        $n=\suc{j}$ by \obv{natural_zero_or_successor}.
    \begin{byCase}
        \caseOf{$n=\emptyset$.}
            $k+n=k$ by \auto{num_addition_zero_right}.
        \caseOf{$n\neq\emptyset$.}
            Take $j$ such that $j$ is a natural number and
                $n=\suc{j}$.
            $j\in n$ by \obv{suc_intro_self}.
            Take $w\in\surrealIntegers$ such that
                $\surrealMul{\surrealEpsilon{k}}{z}
                    =\surrealMul{\surrealEpsilon{k+j}}{w}$.
            Let $w'=\surrealAdd{w}{w}$.
            $w'\in\surrealIntegers$
                by \auto{surreal_integer_addition_closed}.
            $k+j$ and $k+n$ are natural numbers
                by \obv{num_addition_on_naturals}.
            $k+n=\suc{k+j}$
                by \auto{num_addition_successor_right}.
            $\surreal{w}$ and
                $\surreal{\surrealEpsilon{\suc{k+j}}}$
                by \auto{surreal_integer_is_surreal,surreal_epsilon_is_surreal}.
            $\surrealAdd{\surrealEpsilon{\suc{k+j}}}
                {\surrealEpsilon{\suc{k+j}}}
                =\surrealEpsilon{k+j}$
                by \obv{surreal_epsilon_halving,natural_number_is_ordinal}.
            $\surrealMul{\surrealEpsilon{\suc{k+j}}}{w'}
                =\surrealAdd
                    {\surrealMul{\surrealEpsilon{\suc{k+j}}}{w}}
                    {\surrealMul{\surrealEpsilon{\suc{k+j}}}{w}}$
                by \auto{surreal_mul_distributive_left}.
            Hence
                $\surrealMul{\surrealEpsilon{\suc{k+j}}}{w'}
                =\surrealMul
                    {\surrealAdd{\surrealEpsilon{\suc{k+j}}}
                        {\surrealEpsilon{\suc{k+j}}}}{w}$
                by \obv{surreal_mul_distributive_right}.
            Hence
                $\surrealMul{\surrealEpsilon{\suc{k+j}}}{w'}
                =\surrealMul{\surrealEpsilon{k+j}}{w}$.
            Follows by \auto.
    \end{byCase}
\end{proof}

\begin{lemma}\label{surreal_mul_four_rearrange}
    Suppose $\surreal{a}$ and $\surreal{b}$ and
        $\surreal{c}$ and $\surreal{d}$.
    Then
        $\surrealMul{\surrealMul{a}{b}}{\surrealMul{c}{d}}
        =\surrealMul{\surrealMul{a}{c}}{\surrealMul{b}{d}}$.
\end{lemma}
\begin{proof}
    Follows by
        \auto{surreal_mul_associative,surreal_mul_commutative,surreal_mul_is_surreal}.
\end{proof}

\subsection{The dyadic carrier}

A dyadic surreal is an integer multiple of one of the canonical negative
powers of two.  The paired replacement keeps the carrier visibly set-sized.

\begin{definition}\label{surreal_dyadics}
    $\surrealDyadics=
        \{\surrealMul{\surrealEpsilon{\fst{p}}}{\snd{p}}
            \mid p\in\naturals\times\surrealIntegers\}$.
\end{definition}

\begin{proposition}\label{surreal_dyadic_iff}
    $x\in\surrealDyadics$ iff there exist $k,z$ such that
        $k\in\naturals$ and $z\in\surrealIntegers$ and
        $x=\surrealMul{\surrealEpsilon{k}}{z}$.
\end{proposition}
\begin{proof}
    Follows by
        \auto{surreal_dyadics,times_elem_is_tuple,times_tuple_intro,fst_eq,snd_eq}.
\end{proof}

\begin{proposition}\label{surreal_dyadic_is_surreal}
    For every $x\in\surrealDyadics$ we have $\surreal{x}$.
\end{proposition}
\begin{proof}
    Fix $x\in\surrealDyadics$.
    Take $k,z$ such that $k\in\naturals$ and
        $z\in\surrealIntegers$ and
        $x=\surrealMul{\surrealEpsilon{k}}{z}$
        by \auto{surreal_dyadic_iff}.
    $\surreal{\surrealEpsilon{k}}$
        by \auto{surreal_epsilon_is_surreal}.
    $\surreal{z}$ by \obv{surreal_integer_is_surreal}.
    Follows by \obv{surreal_mul_is_surreal}.
\end{proof}

\begin{theorem}\label{surreal_integer_is_dyadic}
    Let $z\in\surrealIntegers$.
    Then $z\in\surrealDyadics$.
\end{theorem}
\begin{proof}
    $\emptyset\in\naturals$ by \obv{omega_empty}.
    $\surreal{z}$ by \obv{surreal_integer_is_surreal}.
    $\surrealMul{\surrealEpsilon{\emptyset}}{z}=z$
        by \auto{surreal_epsilon_zero,surreal_mul_one_left}.
    Follows by \auto{surreal_dyadic_iff}.
\end{proof}

\begin{theorem}\label{surreal_epsilon_is_dyadic}
    Let $k$ be a natural number.
    Then $\surrealEpsilon{k}\in\surrealDyadics$.
\end{theorem}
\begin{proof}
    $\suc{\emptyset}$ is a natural number
        by \obv{omega_empty,omega_successor}.
    Hence $\surrealOne\in\surrealIntegers$
        by \obv{natural_surreal_is_integer}.
    $\surrealMul{\surrealEpsilon{k}}{\surrealOne}
        =\surrealEpsilon{k}$
        by \auto{surreal_mul_one_right,surreal_epsilon_is_surreal}.
    Follows by \auto{surreal_dyadic_iff}.
\end{proof}

\begin{theorem}\label{surreal_dyadic_negation_closed}
    Let $x\in\surrealDyadics$.
    Then $\surrealNeg{x}\in\surrealDyadics$.
\end{theorem}
\begin{proof}
    Take $k,z$ such that $k\in\naturals$ and
        $z\in\surrealIntegers$ and
        $x=\surrealMul{\surrealEpsilon{k}}{z}$
        by \auto{surreal_dyadic_iff}.
    $\surrealNeg{z}\in\surrealIntegers$
        by \obv{surreal_integer_negation_closed}.
    $\surreal{\surrealEpsilon{k}}$ and $\surreal{z}$
        by \auto{surreal_epsilon_is_surreal,surreal_integer_is_surreal}.
    $\surrealNeg{x}
        =\surrealMul{\surrealEpsilon{k}}{\surrealNeg{z}}$
        by \auto{surreal_mul_neg_right}.
    Follows by \auto{surreal_dyadic_iff}.
\end{proof}

\begin{theorem}\label{surreal_dyadic_addition_closed}
    Let $x,y\in\surrealDyadics$.
    Then $\surrealAdd{x}{y}\in\surrealDyadics$.
\end{theorem}
\begin{proof}
    Take $k,z$ such that $k\in\naturals$ and
        $z\in\surrealIntegers$ and
        $x=\surrealMul{\surrealEpsilon{k}}{z}$
        by \auto{surreal_dyadic_iff}.
    Take $l,w$ such that $l\in\naturals$ and
        $w\in\surrealIntegers$ and
        $y=\surrealMul{\surrealEpsilon{l}}{w}$
        by \auto{surreal_dyadic_iff}.
    $k$ and $l$ are ordinals by \obv{natural_number_is_ordinal}.
    Take $a\in\surrealIntegers$ such that
        $\surrealMul{\surrealEpsilon{k}}{z}
            =\surrealMul{\surrealEpsilon{k+l}}{a}$
        by \obv{surreal_epsilon_integer_rescale}.
    Take $b\in\surrealIntegers$ such that
        $\surrealMul{\surrealEpsilon{l}}{w}
            =\surrealMul{\surrealEpsilon{l+k}}{b}$
        by \obv{surreal_epsilon_integer_rescale}.
    $k+l=l+k$ by \obv{natural_addition_commutative}.
    Hence $y=\surrealMul{\surrealEpsilon{k+l}}{b}$
        by \auto.
    Hence $x=\surrealMul{\surrealEpsilon{k+l}}{a}$.
    Let $c=\surrealAdd{a}{b}$.
    $c\in\surrealIntegers$
        by \auto{surreal_integer_addition_closed}.
    $k+l$ is a natural number by \obv{num_addition_on_naturals}.
    $\surrealAdd{x}{y}=\surrealMul{\surrealEpsilon{k+l}}{c}$
        by \auto{surreal_mul_distributive_left,surreal_epsilon_is_surreal,surreal_integer_is_surreal}.
    Follows by \auto{surreal_dyadic_iff}.
\end{proof}

\begin{theorem}\label{surreal_dyadic_multiplication_closed}
    Let $x,y\in\surrealDyadics$.
    Then $\surrealMul{x}{y}\in\surrealDyadics$.
\end{theorem}
\begin{proof}
    Take $k,z$ such that $k\in\naturals$ and
        $z\in\surrealIntegers$ and
        $x=\surrealMul{\surrealEpsilon{k}}{z}$
        by \auto{surreal_dyadic_iff}.
    Take $l,w$ such that $l\in\naturals$ and
        $w\in\surrealIntegers$ and
        $y=\surrealMul{\surrealEpsilon{l}}{w}$
        by \auto{surreal_dyadic_iff}.
    Let $c=\surrealMul{z}{w}$.
    $c\in\surrealIntegers$
        by \auto{surreal_integer_multiplication_closed}.
    $k+l$ is a natural number by \obv{num_addition_on_naturals}.
    $\surreal{\surrealEpsilon{k}}$ and
        $\surreal{\surrealEpsilon{l}}$
        by \auto{surreal_epsilon_is_surreal}.
    $\surreal{z}$ and $\surreal{w}$
        by \obv{surreal_integer_is_surreal}.
    $\surrealMul{x}{y}
        =\surrealMul
            {\surrealMul{\surrealEpsilon{k}}{\surrealEpsilon{l}}}
            {c}$
        by \auto{surreal_mul_four_rearrange}.
    Hence $\surrealMul{x}{y}
        =\surrealMul{\surrealEpsilon{k+l}}{c}$
        by \auto{surreal_epsilon_product,natural_number_is_ordinal}.
    Follows by \auto{surreal_dyadic_iff}.
\end{proof}
